\subsection{A supersingular logarithm estimate excludes the remaining level $81$}

\begin{lemma}[A coefficientwise estimate for a product-type special fibre]
\label{lem:B-RD34-log}
Let $R$ be a complete discrete valuation ring of mixed characteristic
$(0,3)$, with perfect residue field, uniformizer $\pi$, and normalized
valuation $v(3)=a$. Let $G$ be a commutative formal group of finite height
and dimension $d$ over $R$. Suppose the coordinates of $[3]_{G_k}$ have
no monomial of total degree less than nine. Define
\[
 q_0=1,\qquad q_{i+1}=\min\{a+q_i,\,1+3q_i,\,9q_i\}.
                                                        \tag{B-RD34.1}
\]
If $q_u>a/2$, then for every $h\in tR[[t]]^d$ the normalized logarithm
satisfies
\[
       3^u\log_G(\pi h)\in\pi^{q_u}R[[t]]^d.
                                                        \tag{B-RD34.2}
\]
The hypothesis holds for a special fibre isomorphic to a product of
one-dimensional height-two formal groups, and is invariant under an
arbitrary integral change of formal coordinates.
\end{lemma}
\begin{proof}
Write the matrix of normalized invariant differentials as $U(X)$;
its entries are in $R[[X]]$ and $U(0)=1$. Multiplication by three gives
the exact differential identity
\[
     U([3]_G X)\,J_{[3]_G}(X)=3U(X).
                                                        \tag{B-RD34.3}
\]
Since $U([3]_GX)$ is invertible over $R[[X]]$, every entry of the
Jacobian $J_{[3]_G}$ is divisible by $3$. If a monomial $X^\alpha$ has
an exponent not divisible by three, its coefficient in $[3]_G$ is
therefore divisible by three. In particular this holds for every
monomial of total degree two. The linear part is $3X$. The special-fibre
hypothesis makes every coefficient of total degree three through eight
divisible by $\pi$, and all higher coefficients are integral. Thus an
augmentation vector whose coefficients have valuation at least $q\geq1$
is sent by $[3]_G$ to one whose coefficients have valuation at least
\[
                        \min\{a+q,1+3q,9q\}.
\]
Iteration starting with $\pi h$ proves the bound $q_i$ for
$[3]^i_G(\pi h)$.

For completeness this also controls the multivariable logarithm.
Write a coefficient of $X^\alpha$ in any coordinate of $\log_G$ as
$c_\alpha$, and put $m=|\alpha|\geq1$. Since $d\log_G=U(X)dX$,
\[
 v(c_\alpha)\geq
       -a\min_{j:\alpha_j>0}v_3(\alpha_j)
       \geq-a\lfloor\log_3m\rfloor.
                                                        \tag{B-RD34.4}
\]
For every integer $r\geq0$, $3^r-1\geq2r$. Consequently
$\lfloor\log_3m\rfloor\leq(m-1)/2$. If a vector $z$ has coefficient
valuations at least $q>a/2$, every coefficient of
$c_\alpha z^\alpha$ has valuation at least
$mq-a\lfloor\log_3m\rfloor\geq q$. Formal substitution is legitimate
because $z$ is an augmentation vector; each coefficient involves only
finitely many terms. Hence $\log_G(z)\in\pi^qR[[t]]^d$.
Apply this to $z=[3]^u_G(\pi h)$ and use
$\log_G([3]^u_G(\pi h))=3^u\log_G(\pi h)$.

For a one-dimensional height-two group in characteristic three,
multiplication by three starts in degree nine. This persists under
products. An integral coordinate change and its inverse have zero
constant terms and invertible linear parts. Substitution in either
direction preserves the ideal generated by monomials of total degree
at least nine, so the hypothesis is coordinate independent. No
coordinate choice on the mixed-characteristic quotient is assumed.
\end{proof}

\begin{lemma}[A refined kernel estimate]
\label{lem:B-RD34-kernel}
Keep $R,a$ as above. Let $f:F\to G$ be a homomorphism of finite-height
formal groups having no nonzero finite flat subgroup scheme in its
kernel, and with finite special-fibre kernel. This is absence of a
flat kernel; it does not require finite-level closed immersions and
does not preclude a nonzero kernel after reduction. Suppose $G_k$
satisfies the hypothesis of
Lemma~\ref{lem:B-RD34-log}. If $q_u>a/2$, then
\[
                         F^u\ker(f_k)=0,
                                                        \tag{B-RD34.5}
\]
where $F^u$ denotes iterated relative Frobenius. In particular one can
take $u=2$ when $a\leq24$, and $u=3$ when $a=72$.
\end{lemma}
\begin{proof}
We give the change to the original Cartier-module proof explicitly.
Use Bondarko's invariant module $D_G\subset L[[\Delta]]^d$, where
$L=\operatorname{Frac}(R)$, and put
\[
 \mathcal R=R[[\Delta]][1/3],\qquad
 J=\pi^{\lceil a/2\rceil}R.
\]
Here $\mathcal R$ is precisely the ring of series with coefficient
valuations bounded below. Let $G_\pi$ be the formal group law
$\pi^{-1}G(\pi X,\pi Y)$, and let $\mathcal C_G$ denote Bondarko's
completed module $M_G$. We use the exact identities
\[
 \mathcal C_G=\Omega\,\pi D_{G_\pi},\qquad
 J[[\Delta]]^d\subset D_G,\qquad
 E(G_k)=D_G/(D_G\cap\mathcal R^d).
                                                        \tag{B-RD34.6}
\]
They are respectively Proposition~10(4), Proposition~9, and Lemma~3
of \cite[printed pp.~674--675]{Bondarko2006Generic}. The canonical
embedding $D_G\to\mathcal C_G\subset L\{\{\Delta\}\}^d$ agrees
with the literal series on every bounded coefficient projection;
this is Proposition~11(1)--(2), printed p.~675.

By the definition of $D_{G_\pi}$, an element of $\pi D_{G_\pi}$,
viewed as a $3$-typical series in $t$, has the form
$\log_G(\pi h)$ for $h\in tR[[t]]^d$.
Lemma~\ref{lem:B-RD34-log} therefore gives
\[
             3^u\mathcal C_G
                  \subset\pi^{q_u}R\{\{\Delta\}\}^d
                  \subset J\{\{\Delta\}\}^d.
                                                        \tag{B-RD34.7}
\]
Indeed $\Omega$ consists of two-sided series with integral Witt
coefficients, tending to zero towards negative powers. Its action
preserves every coefficient valuation ideal; the formula for this
action is displayed immediately before Proposition~10 in the cited
source. Thus the estimate on $\pi D_{G_\pi}$ passes to its
$\Omega$-span. This replaces the general logarithm bound in
Proposition~11(5).

Here is the needed replacement for Proposition~12(1)(ii). If
$z\in D_G$, $A$ is an integral matrix, and $Az\in\mathcal R^b$, then
\[
                     \mathbf f^uAz\in J[[\Delta]]^b.
                                                        \tag{B-RD34.8}
\]
To see this, (B-RD34.7) and the bounded-projection property give
$3^uAz\in J[[\Delta]]^b$. The Cartier Frobenius satisfies
\[
        \mathbf fz=3(z-z(0))\Delta^{-1},\qquad z(0)\in R^d.
\]
Each $\mathbf f^jz$ again lies in $D_G$, so its constant term is
integral. Induction gives
\[
   \mathbf f^uAz\equiv3^uAz\Delta^{-u}
                        \pmod{3R((\Delta))^b}.
\]
Since $3R\subset J$, the right side and the error lie in
$J((\Delta))^b$. The left side has no negative powers, proving
(B-RD34.8). This is equation~(8) and the proof of Proposition~12,
printed p.~676, with the improved bound inserted.

We now spell out the final step of Proposition~16, printed p.~681.
Let $B=df(0)$, $I'=B(L^{\dim F})$, and $I=I'\cap R^d$.
Choose an integral direct-sum complement $R^d=I\oplus C$ and
the associated projection $\operatorname{pr}:R^d\to C$.
Propositions~6 and~8 identify
\[
                     B D_F=D_G\cap I'[[\Delta]].
                                                        \tag{B-RD34.9}
\]
Indeed the formal group belonging to the right side factors $f$;
the resulting map from $F$ is an isogeny, and generic injectivity
forces its finite flat kernel to vanish, so it is an isomorphism.
The same reduction identity used in Proposition~16 then embeds
$E(F_k)$ in $E(G_k)$.

Take $z\in D_G$ whose reduction $\bar z$ belongs to
$\mathbf Q_3E(F_k)\cap E(G_k)$. By (B-RD34.6), after multiplying
by a power of three, $z$ differs from an element of $BD_F$ by a
bounded-coefficient series. Hence
$\operatorname{pr}(z)\in C\otimes_R\mathcal R$.
Formula (B-RD34.8) shows
\[
 c=\mathbf f^u\operatorname{pr}(z)\in J C[[\Delta]]
                                                \subset D_G.
\]
This element is bounded, so reduces to zero in $E(G_k)$. Moreover
$\mathbf f^uz-c$ lies in the right side of (B-RD34.9). Therefore
$\mathbf f^u\bar z\in E(F_k)$. The saturation quotient
\[
 (\mathbf Q_3E(F_k)\cap E(G_k))/E(F_k)
\]
is the Dieudonne module of $\ker(f_k)$ by Theorem~2(8), equation~(6),
printed p.~668 of the same source. The operator $\mathbf f$
corresponds to relative Frobenius, as stated immediately after that
equation. Full faithfulness proves (B-RD34.5).

For the numerical assertions, if $a\leq2$ then $q_2=2a+1$;
if $3\leq a\leq9$ then $q_2=a+4$; and if $9\leq a\leq24$
then $q_2=13$. Each is greater than $a/2$. At $a=72$ the successive
values are $1,4,13,40$, and $40>36$.
\end{proof}

\begin{theorem}[B-RD34: the height-two least index has only two possibilities]
\label{thm:B-RD34}
\status{Proved. Independent parent and obstruction/crystal proof and original-source
audits passed; exact arithmetic checks also passed.}
For the character question B.4 with $p=3$, absolute ramification
$e_K=2$, $\mu_3\subset K$, and height two,
\[
                         2\leq\nu_{\min}(3,2,2)\leq3.
                                                        \tag{B-RD34.10}
\]
Thus an isomorphism of $81$-torsion always determines the requested
prime-to-$3$ character. A failure at level $27$, if one exists,
must compare two whole potentially supersingular groups with
norm-one finite special inertia and different tame quotient orders;
the order of their joint faithful good-reduction image is $18$ or
$36$.
\end{theorem}
\begin{proof}
The lower bound is the explicit level-nine comparison in
Proposition~\ref{prop:B-RD29-p3-level-two}. The reductions in
Theorem~\ref{thm:B-RD31} and Proposition~\ref{prop:B-RD32} show
that only two whole potentially supersingular groups need be
considered. At level $81$, their joint image would have order
$36$, so their common good-reduction DVR $R$ has $a=72$.
After completed unramified base change, form the graph closure
$H_n\subset P=\mathcal M\times\mathcal N$ at level $3^n$.
Proposition~\ref{prop:B-RD33} gives
\[
       (H_n)_k=\ker F_{M_0}^n\times\ker F_{N_0}^n.
                                                        \tag{B-RD34.11}
\]
The quotient $Q=P/H_n$ is a $3$-divisible group over $R$.
Its special fibre is the product of the corresponding Frobenius
twists of $M_0,N_0$, because quotient formation commutes with base
change. In particular it is a product of one-dimensional
height-two formal groups. Use the homomorphism of formal groups
\[
 g:P\longrightarrow Q\times\mathcal M,
                 \qquad z\longmapsto(qz,\operatorname{pr}_{\mathcal M}z).
                                                        \tag{B-RD34.12}
\]
The generic graph of an isomorphism meets $0\times N$ trivially,
so $g$ has no nonzero finite flat subgroup scheme in its kernel.
By (B-RD34.11), its special kernel is
$0\times\ker F_{N_0}^{n}$. Its target special fibre is a product
of three height-two groups, so Lemma~\ref{lem:B-RD34-kernel}
applies. For $n=4,a=72$ it says that $F^3$ kills
$\ker F_{N_0}^4$. This is impossible: on a one-dimensional formal
group that Frobenius image has rank three. This excludes every
remaining level-$81$ case and proves the upper bound.

For a hypothetical level-$27$ comparison, the same construction
with $n=3$ and $a\leq24$ would force $F^2\ker F_{N_0}^3=0$,
again contradicting its rank-three image. The norm-one joint-image
list in (B-RD32.4) therefore leaves only $18$ and $36$. No
existence claim for a remaining pair, and no decision between
least indices two and three, is made by this intermediate bound.
The final decision $\nu_{\min}(3,2,2)=2$ is proved in
Theorem~\ref{thm:B-RD42}.
\end{proof}

\begin{remark}[Exact arithmetic certificate]
The retained program \path{code/verify_supersingular_log_bound.py}
checks $282$ multivariable monomials, $18\,432$ coefficient-valuation
cases, and $40\,000$ logarithm-degree cases. It also verifies every
ramification value $1\leq a\leq24$, the sequence $(1,4,13,40)$ at
$a=72$, the remaining joint orders $18,36$, and both rank-three
Frobenius images. Its exact author run passed in
$0.05011515400110511$ seconds. These are arithmetic checks of the
proof; no existence of a remaining level-$27$ comparison is asserted.
\end{remark}
